\chapter{Introduction}
\label{sec:introduction}

Whether investing in artificial intelligence (AI) makes a firm more competitive is, in
Kemp's words, a theoretical puzzle \parencite[p.~618]{kemp2024competitive}\footnote{[Kemp
(2024), p.~618] ``The prospect of using artificial intelligence (AI) to establish
competitive advantages presents a theoretical puzzle.''}. The promise is large:
intelligent algorithms are expected to automate a substantial share of work and to open
new products and product-development processes. The same literature lists strategic
obstacles: AI can be myopic and hard for managers to steer, and, built on explicit
knowledge, it resembles a general-purpose
technology \parencite[p.~618]{kemp2024competitive}\footnote{[Kemp (2024), p.~618]
``Despite this promise, a growing body of research has highlighted that AI may present
substantial strategic obstacles. AI may be myopic (Balasubramanian, Ye, \& Xu, 2022),
incapable of perceiving interdependencies within a firm (Raisch \& Krakowski, 2021), and
recalcitrant to managerial control (Murray, Rhymer, \& Sirmon, 2021).'' / ``In addition,
AI is a form of explicit knowledge (Broussard, 2018; Shrestha, He, Puranam, \& von
Krogh, 2021), and resembles a general-purpose technology (Teece, 2018).''}. A
technology every competitor can buy is a poor candidate for an advantage. What firms
get out of buying it is an empirical question.

The empirical evidence has grown quickly and does not agree with itself. On German
panel data, AI use raises productivity for the yes/no measure and the intensity measure
alike \parencite[pp.~189, 201]{czarnitzki2023artificial}\footnote{[Czarnitzki et al.
(2023), p.~189] ``We analyze cross-section as well as panel data from the German part of
the European Commission's Community Innovation Survey (CIS).'' / [p.~201] ``We found
that employing AI technologies has a positive and significant impact
on firm productivity. In particular, we showed that both the use of AI and the intensity
with which firms exploit the potential of AI significantly increase both sales and
value-added.''}. Growth from AI investments concentrates among
the larger firms, and industry concentration rises
alongside \parencite[p.~1]{babina2024artificial}\footnote{[Babina et al. (2024),
p.~1] ``AI-powered growth concentrates among larger firms and is associated with higher
industry concentration.''}. Investors, on average, react to AI investment announcements
with a drop of 1.77 percent on the announcement day, and firms with weak IT bases or
poor credit ratings lose more \parencite[p.~373]{lui2022impact}\footnote{[Lui et al.
(2022), p.~373] ``The stock prices of the firms decrease by 1.77\% on the day of the
announcement. Nonmanufacturing firms and firms with weak information technology
capabilities or low credit ratings suffer a more negative impact compared with other
firms.''}. And where adoption is observed over time, it destroys old capabilities while
creating new ones at the human-machine
boundary \parencite[p.~1425]{krakowski2023artificial}\footnote{[Krakowski et al.
(2023), p.~1425] ``AI adoption triggers interrelated substitution and complementation
dynamics, which make humans' traditional competitive capabilities obsolete, while
creating new sources of persistent heterogeneity when humans interact with chess
engines.''}. Positive, negative, and conditional findings stand side by side, often for
the same outcome. The disagreement itself points at the question worth asking: under which
circumstances do the investments pay off.

That is the question this thesis asks: \emph{under what conditions do firm-level AI
investments translate into competitive advantage and superior firm performance?} The
two outcome constructs are kept analytically separate throughout, because a firm can
improve its performance without holding anything a rival cannot copy; the theoretical
grounds for that separation are laid out in Chapter~\ref{sec:background}.

The thesis answers the question with a systematic literature review of the empirical
evidence, searched over the publication window 2015 to 2026. A documented Scopus query identified 432
records, which a journal-quality filter, title-abstract screening, and full-text
screening reduced to a final corpus of 67 empirical studies; every study was coded
along a fixed scheme by two independent coders, and the resulting coding table is the
review's core deliverable. Chapter~\ref{sec:background} condenses the four theoretical
lenses used to interpret the evidence: general-purpose technologies, the resource-based
view, automation and augmentation, and situated AI. Chapter~\ref{sec:method} documents
the search, screening, and coding protocol. Chapter~\ref{sec:results} presents the
findings in two parts, a descriptive mapping of the evidence base and a synthesis of
the conditions the studies identify, ordered into six groups.
Chapter~\ref{sec:discussion} interprets the findings against the theoretical lenses and
names the limits of what the reviewed literature can show, and
Chapter~\ref{sec:conclusion} concludes.
